% ---------------------------------------------------------------------------
%  09-novelty.tex -- the blunt separation of new / ported / engineering
% ---------------------------------------------------------------------------
\section{What is new, what is ported, what is engineering}
\label{sec:novelty}

This section exists so that the paper never has to be vague about provenance.
A TACAS reviewer's first question is ``which part of this is a contribution?'',
and a submission that answers it before being asked is in a much stronger
position than one that lets the reviewer draw the line.

\subsection{New, and proved here}

\begin{enumerate}[label=\textbf{N\arabic*.}, leftmargin=2.2em, itemsep=4pt]

\item \textbf{Concurrent antichains via coverage monotonicity.}
      \cref{thm:coverage} and \cref{cor:probe}.  A subsumption-based passed set
      is mutable --- entries are deleted and replaced --- which normally forbids
      reading it concurrently with writers.  We identify the structural
      property that makes it safe (\emph{a stored zone is retired only into
      something that contains it}), prove that the covered region is therefore
      monotone, and use it to run the dominant $83\%$ of insertion attempts
      under a shared lock.  This is the load-bearing contribution: it is what
      lets \ltsmin{}'s shared-table shape carry an antichain instead of a
      visited set.

\item \textbf{Exactness of the convex-union merge, and its role as a
      concurrency mechanism.}  \cref{thm:hull} shows the facet test decides
      convexity of the union exactly.  The new observation is what that buys in
      a parallel setting: it holds antichains at $\approx 1$ zone per discrete
      state, which is what makes the critical section $O(1)$.  Merging is
      usually presented as a state-count optimisation; here it is a
      \emph{scalability} mechanism, and the state-count win against \uppaal{}
      on \code{phil} is a side effect.

\item \textbf{Adaptive choice between the two finiteness mechanisms.}
      \cref{thm:adapt}.  Pruning (exact zones $+$ $\aLU$) and widening
      ($\extraLU$) are not interchangeable, and which one wins is a property of
      the model.  The engine measures the distinguishing signal --- zones per
      discrete state --- for free and switches once, one-way, at a level
      boundary.  We prove that the resulting \emph{mixed} passed set is sound
      and that the search still terminates, which is not immediate: the
      termination argument has to combine a finite prefix with a finite image.

\item \textbf{Verdict determinism under a nondeterministic merge order.}
      \cref{thm:determinism}.  The stored antichain depends on the thread
      schedule; the answer does not.  Stating this explicitly matters because
      the tool visibly reports different state counts on repeated runs.

\item \textbf{Recycling reclamation for an interned zone store, with its
      soundness condition.}  \cref{prop:gc,prop:gc-cache}.  Recycling record
      slots and indices in place (no renumbering) keeps the passed set's
      $4$-byte references valid across a GC, but it interacts with a
      byte-validated per-thread cache in a way that is a \emph{soundness} bug,
      not a performance bug.  We prove cache invalidation is both necessary and
      sufficient.

\item \textbf{Removing the last serial phase by bucketed collection.}
      \cref{lem:bucket}.  Small, but it is the difference between a level loop
      that scales and one whose serial fraction grows with level width.

\item \textbf{Minimal-graph inclusion as the enabler of a cheap critical
      section.}  The test itself is UDBM's and we claim no novelty for it
      (\cref{lem:mingraph} is stated and proved for self-containedness).  What
      is new is the observation that being able to decide inclusion
      \emph{without decompressing} is precisely what allows the read-side probe
      of \textbf{N1} to be cheap enough to be worth doing --- otherwise the
      shared lock would be held across an $O(\dim^3)$ closure and the two-phase
      protocol would lose to a plain mutex.
\end{enumerate}

\subsection{Ported or standard, and credited}

\begin{itemize}[leftmargin=1.2em, itemsep=3pt]
\item DBM zone algebra and canonicalisation --- classical \cite{Dill89}.
\item $\extraLU$ widening and LU bounds ---
      \cite{BBLP06}; \uppaal{}'s default.
\item The $\aLU$ subsumption test --- \cite{HSW12}; the implementation is a
      port of \tchecker{}'s \code{is\_alu\_le}, brute-force validated against the
      point-wise definition.
\item Unsoundness of $LU$ abstraction under diagonal constraints, and the
      diagonal-preserving fallback --- \cite{Bouyer04}.
\item Minimal-graph zone storage and the $O(|G|)$ inclusion test --- \uppaal{}'s
      UDBM \code{mingraph}; \file{mindbm.hpp} is a re-expression of it in
      \tool{}'s bound encoding, not a link against UDBM.
\item Active-clock reduction --- \cite{DawsYovine96}.
\item Shared passed table with fine-grained locking and no static partition ---
      \ltsmin{} \cite{LTSmin,OpaalLTSmin}.
\item Level-synchronous parallel BFS --- standard; the choice to keep it rather
      than move to work-stealing DFS is ours and is justified by state counts,
      not by novelty.
\end{itemize}

\subsection{Engineering: real, necessary, not a theoretical claim}

These carry no theorem and should be presented as implementation work.  They
are nonetheless where most of the measured speedup lives, and a paper that
omits them cannot explain its own numbers.

\begin{itemize}[leftmargin=1.2em, itemsep=2pt]
\item $64$-way interner sharding with a high-bit shard selector
      (\cref{prop:shard} is elementary but was learned the hard way).
\item \code{SpinRW}: an $8$-byte spin-then-yield reader--writer lock, sized so
      $16$k of them cost less than $16$k \code{std::shared\_mutex}es.
\item Per-thread pooled allocator for zones and discrete vectors, removing the
      process-global \code{malloc} lock from the successor path.
\item Slot-reusing successor generation (zero allocations per successor after
      warm-up).
\item Frontier as $(\dk,\mathit{zidx})$ pairs; discrete state decoded from a
      packed $63$-bit key.
\item Dense mirror for chains past $64$ live zones.
\item Per-stripe compaction with an amortised trigger.
\item Register-local counters instead of per-thread array slots, to avoid false
      sharing on a per-successor increment.
\end{itemize}

\subsection{One-paragraph version, for the abstract}

\begin{quote}\itshape
Zone-based reachability is usually parallelised by sharing a visited set, but a
competitive timed model checker does not store a visited set: it stores an
antichain of zones under subsumption, which is mutable, and whose entries are
deleted and replaced as the search proceeds.  We show that the antichain update
rule of a subsumption-and-merge passed set has a monotonicity property --- the
covered region never shrinks --- that makes the read side of insertion safe to
run concurrently with writers.  Building on it, we present an OpenMP engine in
which $83\%$ of insertions are resolved under a shared lock, in which the
exclusive critical section is $O(1)$ because an exact convex-union merge keeps
antichains at one zone per discrete state, and in which the choice between the
two standard finiteness mechanisms is made adaptively at run time with a proved
soundness and termination guarantee.  Against \uppaal{} the engine explores the
same symbolic state space on three benchmark families --- exactly the same
stored state counts --- and $\approx 30\%$ fewer states on a fourth, while
running $5.6$--$13.8\times$ faster on $32$ cores and scaling at $81$--$87\%$
efficiency from $16$ to $32$.
\end{quote}
